% Section 1 — Project overview and objective
\section{Project Overview and Objective}
\label{sec:overview}

\subsection{Problem context}
\label{sec:problem}
Retail inventory decisions balance two costs. Ordering too much ties up cash
in holding cost; ordering too little forfeits margin through stockouts. The
preferred order quantity lies between these extremes, and locating it begins
with demand prediction but does not end there. A forecast is an estimate; an
order commits resources. This study covers the full path from estimate to
committed cost.

Two demand environments frame the analysis. The first, captured by the Walmart
M5 data, is sparse and intermittent: most items sell nothing on most days,
with a 64.5\% zero rate across the studied series. Classical forecasting tools
were developed largely for smooth, continuous demand and are poorly suited to
this setting. The second, a store-item demand dataset, is dense: nearly every
item sells every day (zero rate 0.02\%). A method that performs well in one
environment may perform poorly in the other, so any comparative claim is
tested in both.

\subsection{Why forecast accuracy alone is insufficient}
\label{sec:why-not-accuracy}
Forecast accuracy and inventory cost weight errors differently. Consider two
forecasters where A has lower average error than B. If A's errors concentrate
in under-prediction on high-volume days while B's errors are modest
over-predictions on quiet days, B can produce the lower inventory cost under a
real ordering rule. Accuracy metrics treat each unit of error symmetrically,
while inventory costs penalize stockouts and overstocks asymmetrically, filter
forecasts through safety-stock formulas, and interact with lead times and
review rhythms. Evaluating a forecasting model without passing its forecasts
through an inventory policy therefore measures only part of its economic
effect. Every model in this study is evaluated twice: once as a forecaster
and once as an input to inventory decisions under a single shared ordering
policy.

\subsection{Research question}
\label{sec:rq}
How does the effectiveness of AI-based inventory optimization change as we move
from traditional forecasting models toward neural-network approaches, when
effectiveness is measured at the level of inventory outcomes rather than
forecast errors alone?

\subsection{Hypotheses}
\label{sec:hypotheses}
The study tests four neutrally framed claims so that any outcome is
informative: \textbf{H1,} accuracy differs significantly across model
generations; \textbf{H2,} accuracy differences produce measurable inventory
differences under a common policy; \textbf{H3,} sophistication improves
inventory outcomes monotonically (the claim under test; its rejection would be
the finding, not a failure); \textbf{H4,} patterns hold across both retail
environments. Verdicts appear in \cref{sec:verdicts}.

\subsection{Objectives}
\label{sec:objectives}
\begin{enumerate}
\item Compare twelve forecasting approaches (baselines, exponential
smoothing, ARIMA/SARIMA, Croston-family specialists, and a global LSTM) on
identical footing across two demand environments.
\item Evaluate each model on forecast accuracy and on the inventory cost its
forecasts generate under a shared ordering policy.
\item Distinguish statistical significance from practical importance using
paired tests, effect sizes, and confidence intervals.
\item Distinguish policy-stable conclusions from policy-sensitive ones using a
27-policy sensitivity grid.
\item Publish a frozen, traceable evidence base that supports independent
audit of every reported number.
\end{enumerate}

\subsection{Scope}
\label{sec:scope}
The study covers univariate daily demand forecasting (no prices, promotions, or
calendar features), a 28-day horizon over eight rolling origins, and one
family of inventory policies (daily-review order-up-to, lost sales). The M5
analysis uses a stratified sample of 500 series rather than the full 30,490
catalog. SARIMA runs Store-only and the Croston family M5-only, each for
documented methodological reasons. The planned large-language-model phase was
designed but not executed; it appears only as future work (\cref{sec:llm-future}).

\subsection{Research contributions}
\label{sec:contributions}
The empirical contribution is a measurement of the forecast-decision
relationship in both directions within one controlled design: correspondence
on M5 demand and divergence on Store demand. The methodological contribution
comprises a shared-policy evaluation harness that holds the inventory
treatment constant across models, a leakage-safe rolling protocol with
scale-free metrics, and a sensitivity discipline in which headline rankings
are reported together with their stability across policies. The analytical
contribution pairs each significance statement with an effect size and each
headline ranking with a policy-stability assessment. For reproducibility, the
frozen evidence package (number sheet, figure and table audits, claim
traceability) maps every reported number to its source file.

\subsection{Datasets at a glance}
\label{sec:datasets-glance}
\textbf{M5 (primary).} Walmart M5 competition data \citep{makridakis2022m5}:
thousands of item-store daily series, 2011 to 2016. This study uses 500
stratified series over the common window 2013-01-01 to 2016-05-22. The
archetype split is 194 highly intermittent, 199 intermittent, 26 moderate, 80
variable, and 1 smooth series. \textbf{Store Item Demand (secondary robustness
environment).} Kaggle Store Item Demand Forecasting Challenge: 10 stores by 50
items (500 series), a complete panel of 913,000 rows from 2013-01-01 to
2017-12-31, analyzed over the same common window. Demand is dense, with
essentially no zero-demand days. \Cref{fig:rep-m5,fig:rep-store,fig:diversity}
illustrate the contrast.

\begin{figure}[H]
\centering
\includegraphics[width=0.9\textwidth]{04_representative_series.png}
\caption{Representative sparse M5 series: extended zero-demand intervals
punctuated by spikes. Source: 07\_figures/eda/m5 (notebook 03).}
\label{fig:rep-m5}
\end{figure}
\begin{figure}[H]
\centering
\includegraphics[width=0.9\textwidth]{08_representative_series.png}
\caption{Representative dense Store series: nonzero demand each day with a
weekly rhythm. Source: 07\_figures/eda/store\_item\_demand (notebook 03).}
\label{fig:rep-store}
\end{figure}
\begin{figure}[H]
\centering
\includegraphics[width=0.9\textwidth]{11_series_diversity_panel.png}
\caption{Store demand diversity: level and volatility differ widely across
series within a dense environment. Source: notebook 03.}
\label{fig:diversity}
\end{figure}

\subsection{What the study establishes}
\label{sec:establishes}
First, a forecast leaderboard for both datasets with uncertainty quantified.
Second, a decision-level leaderboard showing where forecasting skill reduces
cost and where it does not. Third, a robustness map recording which
conclusions hold across lead times, service targets, and cost ratios, and
which vary with them. In summary: on M5 demand the lowest-error forecaster is
also the lowest-cost inventory input and remains so across policies; on Store
demand it is not, and the lowest-cost input depends on the policy.
